\section{梯形权、积分显式公式与素数提取}\label{sec:trapezoid}

本节把素数存在性归结为三个具有明确含义的非负上界之和小于
主项。论证首先是一个有限和恒等式，然后才使用积分显式公式。
这种顺序使端点、真素数幂和零点截断的处理均可单独核对。
本节使用的显式公式是前述输入
\texttt{IntegratedExplicitFormula}；它本身是经典的无条件恒等式，
其精确来源及与本文约定的对应见文献 \cite{GM2016} 和输入审计。
这里不额外假定 RH，也不把 QRH 用于推导该恒等式。

\subsection{从有限 Mangoldt 和构造梯形权}

对 $t>1$ 定义
\begin{equation}\label{tr:psi1}
 \Psi_1(t)=\sum_{1\leq m\leq\lfloor t\rfloor}\Lambda(m)(t-m).
\end{equation}
若 $t$ 是整数，最后一项为零；因此这一约定没有端点半权问题。
对于任何 $Y\geq t$，也可以写成
\[
 \Psi_1(t)=
 \sum_{1\leq m\leq\lfloor Y\rfloor}\Lambda(m)(t-m)_+,
 \qquad v_+=\max\{v,0\}.
\]
因为新增的每一项均为零。这是下面交换四点差分与有限求和的依据。

设 $x>1$、$h>0$，并设 $0<a\leq b$、$a+b=1$。定义
\begin{align}
 D_{x,h}F
  &=\frac{F(x)-F(x+ah)-F(x+bh)+F(x+h)}{abh},
       \label{tr:difference}\\
 w_{x,h}(t)
  &=\frac{(x-t)_+-(x+ah-t)_+-(x+bh-t)_++(x+h-t)_+}
          {abh}. \label{tr:weight}
\end{align}
在不引起混淆时省略 $x,h$ 下标。逐段展开正部函数得到
\begin{equation}\label{tr:piecewise}
 w(t)=
 \begin{cases}
 0,&t\leq x,\\
 (t-x)/(abh),&x\leq t\leq x+ah,\\
 1/b,&x+ah\leq t\leq x+bh,\\
 (x+h-t)/(abh),&x+bh\leq t\leq x+h,\\
 0,&t\geq x+h.
 \end{cases}
\end{equation}
相邻公式在公共端点取相同值。特别地，
\begin{equation}\label{tr:weight-properties}
 0\leq w(t)\leq \frac1b,\qquad
 w(t)>0\Longrightarrow x<t<x+h,\qquad
 \int_{\R}w(t)\,\dd t=h.
\end{equation}
面积归一化为 $h$，故后面的主项也是 $h$。
本文固定
\begin{equation}\label{tr:fixed-ab}
 a=\frac1{32},\qquad b=\frac{31}{32}.
\end{equation}
该选择不表示在所有参数下都最优；它是当前证书实际采用的权。

\begin{lemma}[有限和与四点差分]\label{tr:finite-identity}
设 $Y=x+h$，并令
\[
 W(x,h)=\sum_{1\leq m\leq\lfloor Y\rfloor}\Lambda(m)w_{x,h}(m).
\]
则 $W(x,h)=D_{x,h}\Psi_1$。此外，
$D_{x,h}(t^2/2)=h$，而 $D_{x,h}$ 消去一切仿射函数。
\end{lemma}
\begin{proof}
将 \eqref{tr:psi1} 的四个截断统一扩展到 $Y$，
并把四个有限和合并，其第 $m$ 项恰为
$\Lambda(m)w_{x,h}(m)$。常数项的系数为 $1-1-1+1=0$；
一次项的系数为 $x-(x+ah)-(x+bh)+(x+h)=0$。
二次项的分子为 $abh^2$，除以 $abh$ 即为 $h$。
\end{proof}

这些有限和及支撑性质对应
\leanref{IntegratedPsi}{weightedMangoldt_trapezoid_eq}、
\leanref{IntegratedPsi}{trapezoidWeight_bounds}
和 \leanref{IntegratedPsi}{trapezoidWeight_pos_support}。

\subsection{零点核的三个互补上界}

对 $\rho=\beta+\ii\gamma\notin\{0,-1\}$ 和 $t>0$，令
\[
 G_\rho(t)=\frac{t^{\rho+1}}{\rho(\rho+1)},\qquad
 K_{x,h}(\rho)=D_{x,h}G_\rho,\qquad
 t^\rho=\exp(\rho\log t).
\]
正实轴上的幂使用实对数，因而这里不存在复对数分支选择。
直接求导得 $G_\rho'(t)=t^\rho/\rho$ 及
$G_\rho''(t)=t^{\rho-1}$。

\begin{lemma}[基本界、一次衰减界与二次衰减界]\label{tr:three-kernels}
在上述条件下，若 $\beta\leq1$，则
\begin{equation}\label{tr:base-kernel}
 |K_{x,h}(\rho)|\leq h x^{\beta-1}.
\end{equation}
若另有 $\gamma\ne0$，则
\begin{align}
 |K_{x,h}(\rho)|
 &\leq h x^{\beta-1}\frac{2}{b}
                  \frac{x+h}{h|\gamma|},\label{tr:first-kernel}\\
 |K_{x,h}(\rho)|
 &\leq h x^{\beta-1}\frac{4}{ab}
                  \frac{(x+h)^2}{h^2\gamma^2}. \label{tr:second-kernel}
\end{align}
\end{lemma}
\begin{proof}
两次微积分基本定理给出
\[
 K_{x,h}(\rho)=
 \frac1{abh}\int_0^{ah}\int_0^{bh}
       (x+u+v)^{\rho-1}\,\dd v\,\dd u.
\]
由于 $\beta-1\leq0$，被积函数的模不超过
$x^{\beta-1}$，得到 \eqref{tr:base-kernel}。
Lean 中使用两次导数范数的增量界证明同一估计。

将分子分成两段长度为 $ah$ 的增量。
在 $x\leq t\leq x+h$ 上，
\[
 |G_\rho'(t)|
 =\frac{t^\beta}{|\rho|}
 \leq \frac{(x+h)x^{\beta-1}}{|\gamma|}.
\]
两个增量的模之和除以 $abh$，得到
\eqref{tr:first-kernel}。最后，
$|\rho|\geq|\gamma|$、$|\rho+1|\geq|\gamma|$，且
$t^{\beta+1}=t^2t^{\beta-1}\leq(x+h)^2x^{\beta-1}$。
分别估计四个函数值即得 \eqref{tr:second-kernel}。
\end{proof}

令
\begin{equation}\label{tr:height-kernel}
 c_1=\frac{64}{31},\qquad c_2=\frac{4096}{31},\qquad
 g_q(t)=\min\left\{1,\frac{c_1q}{t},\frac{c_2q^2}{t^2}\right\}
 \quad(t>0,\ q\geq0).
\end{equation}
由 \eqref{tr:fixed-ab}，$c_1=2/b$，$c_2=4/(ab)$。
若 $q\geq(x+h)/h$，三个界合并为
\begin{equation}\label{tr:kernel-min}
 |K_{x,h}(\rho)|\leq h x^{\beta-1}g_q(|\gamma|).
\end{equation}
$g_q$ 关于高度单调不增，关于 $q$ 单调不减，而且非负。
其两个衰减分支具有不同用途：一次分支改善有限高度范围的估计，
二次分支使无限高度余项可被几何级数控制。
对应的形式化结果为
\leanref{TrapezoidKernel}{norm_trapezoidKernel_le_base}、
\leanref{TrapezoidKernel}{norm_trapezoidKernel_le_first}、
\leanref{TrapezoidKernel}{norm_trapezoidKernel_le_second}
及 \leanref{TrapezoidKernel}{norm_trapezoidKernel_le_min}。

\subsection{积分显式公式的四点作用及极限}

记 $m_\rho$ 为非平凡零点的重数。使用按
$|\operatorname{Im}\rho|\leq T$ 截断的零点和，输入公式可写为
\begin{equation}\label{tr:explicit}
 \Psi_1(t)=\frac{t^2}{2}
 -\lim_{T\to+\infty}
   \sum_{|\operatorname{Im}\rho|\leq T}m_\rho G_\rho(t)
 -\frac{\zeta'}{\zeta}(0)t+\frac{\zeta'}{\zeta}(-1)-Q(t),
 \qquad t>1,
\end{equation}
其中
\begin{equation}\label{tr:trivial-primitive}
 Q(t)=\sum_{j=0}^{\infty}
       \frac{t^{-2j-1}}{2(j+1)(2j+1)}
     =\sum_{j=0}^{\infty}G_{-2(j+1)}(t).
\end{equation}
这里的两个仿射常数是精确常数，不能把未指定的 $O(t)$
直接当作可消去的仿射函数。不过，一旦采用精确公式
\eqref{tr:explicit}，其数值完全不进入后续预算，因为四点差分消去它们。

定义
\[
 Z_T(x,h)=\sum_{|\operatorname{Im}\rho|\leq T}
                  m_\rho K_{x,h}(\rho),\qquad E(x,h)=D_{x,h}Q .
\]
对四个大于 $1$ 的实数同时使用 \eqref{tr:explicit}，并对四个极限作有限线性组合，
由引理 \ref{tr:finite-identity} 得
\begin{equation}\label{tr:zero-limit}
 Z_T(x,h)\longrightarrow h-W(x,h)-E(x,h).
\end{equation}
这一步只需要输入已经给出的高度截断极限；无需在四点操作中
另行假定非平凡零点级数绝对收敛。
因此，只要对每个 $T$ 证明 $|Z_T(x,h)|\leq hZ$，
闭球对极限封闭就给出
\begin{equation}\label{tr:weighted-lower}
 W(x,h)\geq h(1-Z)-|E(x,h)|.
\end{equation}
本节以下所有零点预算都对任意截断高度成立，正是为了这一极限步骤。
对应
\leanref{IntegratedExplicitFormula}{trapezoid_zero_sum_tendsto}。

\begin{lemma}[平凡零点的完整上界]\label{tr:trivial-bound}
若 $x>1$，则
\begin{equation}\label{tr:V}
 |E(x,h)|\leq h V(x),\qquad
 V(x)=\frac{x^{-3}}{1-x^{-2}}.
\end{equation}
\end{lemma}
\begin{proof}
对 $t>1$，\eqref{tr:trivial-primitive} 中第 $j$ 项的模
不超过 $t^{-1}(t^{-2})^j$，故其级数绝对收敛。
对 $\rho=-2(j+1)$ 使用基本核界 \eqref{tr:base-kernel} 得
\[
 |K_{x,h}(-2(j+1))|
 \leq h x^{-2j-3}=h x^{-3}(x^{-2})^j.
\]
此级数也绝对收敛。于是四个收敛级数的有限线性组合可以逐项进行，
且几何级数求和给出 \eqref{tr:V}。
\end{proof}
对应 \leanref{TrapezoidTrivial}{norm_trapezoidTrivialDifference_le}。

\subsection{真素数幂的显式预算}

记 $\theta(Y)=\sum_{p\leq Y}\log p$、
$\psi(Y)=\sum_{m\leq Y}\Lambda(m)$。由 $0\leq w\leq1/b$，
非素数项的总贡献满足
\begin{equation}\label{tr:nonprime}
 \sum_{\substack{1\leq m\leq Y\\m\ \text{非素数}}}
       \Lambda(m)w(m)
 \leq\frac1b\bigl(\psi(Y)-\theta(Y)\bigr)
 \leq\frac{\sqrt Y(\log Y)^2}{2b\log2}.
\end{equation}
后一界可以直接由真素数幂计数看出：
若 $p^\ell\leq Y$ 且 $\ell\geq2$，则
$\ell\leq\log Y/\log2$，每个指数对应的底数数量不超过 $\sqrt Y$，
而 $\log p\leq(\log Y)/2$；对这些可能性求和只会多计。
在本文使用范围内，Lean 另由
$\psi(Y)-\theta(Y)\leq2\sqrt Y\log Y$ 和 $\log Y\geq4\log2$
推出同一个较宽上界。
两种推导都无须假定零点分布，也无须判定区间内有哪些素数幂。
见 \leanref{PrimePowerBudget}{weightedNonprime_le_psi_sub_theta}
及 \leanref{PrimePowerBudget}{weightedNonprime_le_manuscript_budget}。

为了给出完全确定的单元预算，令 $k\in\{4,6\}$，
\begin{equation}\label{tr:domains}
 x=\e^L,\qquad h_k(L)=k\e^{(1-1/k)L},\qquad
 \epsilon_k(A)=k\e^{-A/k},\qquad
 q_k(L)=1+\frac{\e^{L/k}}{k}.
\end{equation}
于是 $x+h_k(L)=\e^L(1+\epsilon_k(L))$，且
$q_k(L)=(x+h_k(L))/h_k(L)$。对 $0<A\leq L\leq C$ 定义
\begin{equation}\label{tr:Pcell}
 P_k(A,C)=
 \frac{32}{31}\,
 \frac{\sqrt{1+\epsilon_k(A)}
       \bigl[C+\log(1+\epsilon_k(A))\bigr]^2}
      {2k\log2}\,
 \e^{-(1/2-1/k)A},
 \qquad
 V_A=\frac{\e^{-3A}}{1-\e^{-2A}}.
\end{equation}
这些是后文 $P$ 和 $V$ 的完整定义。特别是
$P_6$ 的分母为 $12\log2$、指数为 $-A/3$；
$P_4$ 的分母为 $8\log2$、指数为 $-A/4$。

\begin{lemma}[单元内的统一小项界]\label{tr:cell-small}
在本文的中段范围 $167\leq A\leq L\leq C$（$k=6$）
或远段范围 $400\leq A\leq L\leq C$（$k=4$）内，
非素数加权和不超过 $h_k(L)P_k(A,C)$，
且 $|E(\e^L,h_k(L))|\leq h_k(L)V_A$。
\end{lemma}
\begin{proof}
把 \eqref{tr:nonprime} 除以 $h_k(L)$，使用
\[
 \frac{\sqrt{x+h_k(L)}}{h_k(L)}
 =\frac1k\sqrt{1+\epsilon_k(L)}
           \e^{-(1/2-1/k)L}.
\]
函数 $\epsilon_k$ 递减，且 $1/2-1/k>0$。
又有
$\log(x+h_k(L))\leq C+\log(1+\epsilon_k(A))$。
所有因子非负，逐因子放大即得 $P_k(A,C)$。
对平凡零点，$\e^{-3L}\leq\e^{-3A}$，
$1-\e^{-2L}\geq1-\e^{-2A}>0$，使用 \eqref{tr:V} 即可。
\end{proof}
此处分别对应
\leanref{MiddleDomains}{middle_prime_power_normalized_le}、
\leanref{MiddleDomains}{middle_prime_power_le}
和 \leanref{FarDomains}{far_prime_power_le}。

\subsection{严格不等式如何产生真正的素数}

\begin{proposition}[三项预算判据]\label{tr:prime-criterion}
设输入积分显式公式成立，$x>1$、$h>0$。
若对任意实数 $T$ 均有 $|Z_T(x,h)|\leq hZ$，
平凡零点差分满足 $|E(x,h)|\leq hV$，
非素数加权和不超过 $hP$，且
\begin{equation}\label{tr:criterion}
 Z+P+V<1,
\end{equation}
则存在素数 $p$ 满足 $x<p<x+h$。
\end{proposition}
\begin{proof}
由 \eqref{tr:weighted-lower}，
$W(x,h)\geq h(1-Z-V)>hP$。
减去不超过 $hP$ 的非素数项，剩下的素数加权和严格为正。
这是有限个非负项的和，所以至少有一项
$\log p\,w(p)>0$。由 \eqref{tr:weight-properties} 的支撑性质，
$x<p<x+h$。由于两个端点权均为零，结论是严格开区间，
不存在把端点素数误计入结论的问题。
\end{proof}
该结论对应 \leanref{TrapezoidPrime}{prime_of_trapezoid_budgets}。
后续任务因此明确为：构造完整零点预算 $Z$，再与
\eqref{tr:Pcell} 中已定义的 $P_k(A,C)$、$V_A$ 相加。
本节没有把任何待证的素数存在命题作为输入。
